\section{Simulation of Random Variables}

\begin{example}[Simulating of Random Variable from Uniform Distribution]
    \label{ex:simulation-from-uniform}%
    Suppose \(X\) is a random variable with probability distribution function \(F\pqty{x} = \Prob\pqty{X \leq x}\), and suppose we know how to simulate \(U \sim \Uniform \ccintv{0}{1}\). How may we generate \(X\)?

    \begin{itemize}
        \item If \(F\) is one-to-one, i.e., \(F\) is strictly increasing. Set
        \[
            X = F\inverse\pqty{U},
        \]
        then
        \[
            \Prob\pqty{X \leq x} = \Prob\pqty{F\inverse\pqty{U} \leq x} = \Prob\pqty{U \leq F\pqty{x}} = F\pqty{x}
        \]
        gives the desired distribution.

        \item Suppose that \(F\) is not one-to-one. Then we may define its \emph{generalised inverse} as
        \[
            G\pqty{u} = \inf \set{x \in \RR}{u \leq F\pqty{x}}.
        \]

        We claim that \(G\pqty{u} \leq x\) if and only if \(u \leq F\pqty{x}\). Indeed, if \(u \leq F\pqty{x}\), then by definition, \(G\pqty{u} \leq x\). On the other hand, note that there must be some sequence \(x_n \geq G\pqty{u}\),such that \(x_n \to G\pqty{u}\), with \(u \leq F\pqty{x_n}\). By the right-continuity of \(F\), we have
        \[
            \lim_{n \to \infty} F\pqty{x_n} = F\pqty{G\pqty{u}}.
        \]

        So \(u \leq F\pqty{G\pqty{u}}\). Therefore, if \(G\pqty{u} \leq x\), then
        \[
            u \leq F\pqty{G\pqty{u}} \leq F\pqty{x}
        \]
        since \(F\) is increasing.

        Now, set \(X = G\pqty{U}\), then
        \[
            \Prob\pqty{X \leq x} = \Prob\pqty{G\pqty{U} \leq x} = \Prob\pqty{U \leq F\pqty{x}} = F\pqty{x}.
        \]
    \end{itemize} 
\end{example}

\begin{example}[Box--Muller Transform]
    \label{ex:box-muller}%
    We want to simulate \(X, Y \sim \Normal\pqty{0, 1}\) being independent, using two independent \(\Uniform\ccintv{0}{1}\) random variables. However, \autoref{ex:simulation-from-uniform} does not work very well, since the distribution function of the normal distribution does not have a closed form.

    For \(\pqty{X, Y}\), we consider the polar transformation \(X = R \cos \Theta\), \(Y = R \sin \Theta\). Recall that \(R\) has density \(r e^{- \flatfrac{r^2}{2}}\) on \(\oointv{0}{\plusinfty}\), and \(\Theta \sim \Uniform \ccintv{0}{2\pi}\), and they are independent, from \autoref{ex:transformation-cartesian-polar}.

    Let \(U, V\) be independent \(\Uniform\ccintv{0}{1}\)s. Then
    \[
        \Theta = 2 \pi U \sim \Uniform \ccintv{0}{2\pi},
    \]
    and
    \[
        R = \sqrt{- 2 \log V},
    \]
    and we have that
    \[
        \Prob\pqty{R \geq r} = \Prob\pqty{\sqrt{- 2 \log V} \geq r} = \Prob\pqty{V \leq e^{- \flatfrac{r^2}{2}}} = e^{- \flatfrac{r^2}{2}}
    \]
    so \(R\) has the correct density. Then this gives \(X\) and \(Y\) as desired.
\end{example}

\begin{example}[Rejection Sampling]
    \label{ex:rejection-sampling}%
    Suppose we have a random variable \(\vb{X}\), with density
    \[
        f\pqty{\vb{x}} = \frac{\indicator \pqty{\vb{x} \in A}}{\abs{A}}
    \]
    where \(\emptyset \neq A \subseteq \RR^d\), and \(\abs{A} < \infty\) is the volume of \(A\). In other words, \(\vb{X}\) is uniform on the set \(A\) in \(\RR^d\).

    Suppose that \(A\) is a subset of \(\ccintv{0}{1}^d\). Let \(\pqty{U_{k, n} : k = 1, \ldots, d, n \in \NN}\) be independent and identically distributed \(\Uniform \ccintv{0}{1}\). Then
    \[
        \vb{U}_n = \pqty{U_{1, n}, \ldots, U_{d, n}} \sim \Uniform\pqty{\ccintv{0}{1}^d}.
    \]

    Set the random variable
    \[
        N = \min \set{n \in \NN}{\vb{U}_n \in A}
    \]
    and set \(\vb{X} = \vb{U}_N\). To show \(\vb{X}\) is indeed our desired distribution, it suffices to check that
    \[
        \int_{B} f\pqty{\vb{x}} \dd[d]{\vb{x}} = \Prob\pqty{\vb{X} \in B} = \frac{\abs{B \cap A}}{\abs{A}}.
    \]

    We have
    \begin{align*}
        \Prob\pqty{\vb{X} \in B} &= \Prob\pqty{\vb{U}_N \in B} \\
        &= \sum_{n = 1}^{\infty} \Prob\pqty{\vb{U}_n \in N, N = n} \\
        &= \sum_{n = 1}^{\infty} \Prob\pqty{\vb{U}_n \in B, \vb{U}_n \in A, \vb{U}_{n - 1} \notin A, \ldots, \vb{U}_1 \notin A} \\
        &= \sum_{n = 1}^{\infty} \Prob\pqty{\vb{U}_n \in B \cap A} \cdot \pqty{\Prob\pqty{\vb{U}_1 \notin A}}^{n - 1} \\
        &= \sum_{n = 1}^{\infty} \abs{B \cap A} \pqty{1 - \abs{A}}^{n - 1} \\
        &= \frac{\abs{B \cap A}}{\abs{A}}
    \end{align*}
    which gives precisely the desired probability as desired.
\end{example}